%%%%%%%%%%%%%%%%%%%%%%%%%%%%%%%%%%%%%%%%%%%%%%%%%
% EXERCICE
\begin{exercise}[points=3]
    Traduis les propriétés suivantes par des expressions algébriques en n'oubliant pas d'en écrire le résultat.
    
    \begin{center}
        \begin{tblr}{
                hlines,
                vlines,
                rowsep=2ex,
                row{1}={font=\bfseries},
                colspec={XX},
                column{1}={l},
                column{2}={c},
            }
            Nom des propriétés                  & Expressions algébriques\\
            Produit de puissances de même base  & \\
            Puissance d'un produit              & \\
            Puissance d'une puissance           & \\
            Puissance d'un quotient             & \\
            Quotient de puissances de même base & \\
            Puissance dont l'exposant est nul   & 
        \end{tblr}
    \end{center}
\end{exercise}
% SOLUTION
\begin{solution}

\end{solution}
%%%%%%%%%%%%%%%%%%%%%%%%%%%%%%%%%%%%%%%%%%%%%%%%%
%%%%%%%%%%%%%%%%%%%%%%%%%%%%%%%%%%%%%%%%%%%%%%%%%
% EXERCICE

\begin{exercise}[points=2]
    Qu'est-ce qu'une fraction? Sois précis.
    \par\blank[width=3\linewidth,linespread=1.5]{}
\end{exercise}

\pagebreak
\begin{exercise}[points=5]
    Vrai ou faux ? Justifie dans les deux cas.
    \begin{enumerate}
        \item 24 et 33 sont des nombres premiers entre eux.
        \par\blank[width=2\linewidth,linespread=1.5]{}
        \item Une symétrie centrale est une rotation.
        \par\blank[width=2\linewidth,linespread=1.5]{}
        \item $\left(x^2\right)^3 = x^5$.
        \par\blank[width=2\linewidth,linespread=1.5]{}
        \item Si l'amplitude de l'angle d'une rotation est positive alors la figure tourne dans le sens horlogique.
        \par\blank[width=2\linewidth,linespread=1.5]{}
        \item $\dfrac{4}{10}=\dfrac{8}{5}$.
        \par\blank[width=2\linewidth,linespread=1.5]{}
    \end{enumerate}
\end{exercise}
% SOLUTION
\begin{solution}

\end{solution}
%%%%%%%%%%%%%%%%%%%%%%%%%%%%%%%%%%%%%%%%%%%%%%%%%
%%%%%%%%%%%%%%%%%%%%%%%%%%%%%%%%%%%%%%%%%%%%%%%%%
% EXERCICE
\begin{exercise}[points=2]
    Coche la réponse correcte.

    \begin{enumerate}
        \item $2n+1$ et $2n+2$ sont deux nombres \dots
              \begin{qcm}[cols=4]
                  \item pairs
                  \item impairs
                  \item consécutifs
                  \item premiers
              \end{qcm}
        \item Dans l'égalité $49=4\cdot10+9$, le nombre $10$ est \dots
              \begin{qcm}[cols=4]
                  \item le diviseur
                  \item le dividende
                  \item le quotient
                  \item le reste
              \end{qcm}
        \item $(-2)^3$ est égale à \dots
              \begin{qcm}[cols=4]
                  \item $8$
                  \item $6$
                  \item $-8$
                  \item $-6$
              \end{qcm}
        \item $\left(3x^2\right)^2$ se réduit en \dots
              \begin{qcm}[cols=4]
                  \item $3x^4$
                  \item $9x^4$
                  \item $9x^2$
                  \item $6x^4$
              \end{qcm}
    \end{enumerate}
\end{exercise}

\pagebreak
% SOLUTION
\begin{solution}

\end{solution}
%%%%%%%%%%%%%%%%%%%%%%%%%%%%%%%%%%%%%%%%%%%%%%%%%
